\section{El término firmado y su control por cuatro ciclos}
Para $H=H_P$, la resolvente da
$H_D-H=(I-H_D)(O_D-O_P)(I-H)$. Su norma $\HS$ es $O(k^{-1})$ y sus entradas son
$O(k^{-2})$ porque los dos factores externos tienen sumas absolutas por fila y columna acotadas.
Así $V(D)=V(P)+O(k^{-2})$: en los cubos se usa
$\|H\|_{\HS}=O(1)$, en las diagonales $\|\diag_{\rm vec}(H)\|_2=O(k^{-1/2})$,
y en la forma cuadrática la cota de operador del cambio de $H$.
Aquí $\diag_{\rm vec}$ extrae las entradas diagonales.
Sea $v=s-1$. Derivando $\log d_P(y)$ dos veces y usando $\|D-P\|_{\HS}=O(k^{-1})$ se obtiene
\[
 \frac{d_D''}{d_D}=v+v^2+(1+2v)j_0+j_0^2+2\tr T_R^2+O(k^{-1}).
\]
En efecto, para $P$ las dos derivadas de $\log d_P$ en $1$ son
$\tr T_R$ y $\tr T_R+2\tr T_R^2$. Ambas trazas son Lipschitz respecto de
$P$ en norma $\HS$ bajo las cotas uniformes. Esto se ve escribiendo
$G_R(P)-G_R(D)$ como suma de productos de dos operadores $\HS$, de norma de traza $O(k^{-1})$.
En $V$, el vector diagonal es $-(t+j)$ y $H\mathbf1=0$.
La parte de cubos de los bloques diagonales es
$-[\Sigma_3(T_L)+\Sigma_3(T_R)]/3+O(k^{-2})$.
Los cubos cruzados suman $(2/3)\sum(B+U)_{ij}^3$.
Las diagonales omitidas en la forma cuadrática y en los cubos cuestan $O(k^{-2})$.
Para seguir las cancelaciones bastan
\begin{gather*}
 \|t\|^2=Z-4s/k+2/k,\qquad M_z=Z-s^2/k-z^\top H z,\\
 \sum B_{ij}^3=l-3s/k+2/k,\qquad
 C(p\circ U)=j,\quad \langle U,p\rangle=j_0,\quad \mathbf1^\top j=2j_0.
\end{gather*}
Las dos identidades con $U$ se obtienen de sus sumas nulas y de
$UB^\top=T_L-G_L$ y $U^\top B=T_R-G_R$.
Los términos cuadrados en $p$ y $B$ que multiplican $U$ dejan
$2\langle B^{\circ2}-p^{\circ2},U\rangle=-4j_0/k$.
Los términos mixtos restantes con $j$ son
\[
 (\Gamma z)^\top j+z^\top j/2-3t^\top j/2+t^\top Hj=(3-s)j_0/k.
\]
La parte lineal de $d_D''/[2(k-1)d_D]$ aporta $(s-1/2)j_0/k$.
Juntos dan $-3j_0/(2k)$. Los términos libres de $j,T_s,U$ se reducen a
$-1/(12k)$ al sustituir $v=s-1$. Los restantes son los de~\eqref{eq:xi}.
Esta cuenta conserva los signos de los cubos, incluidos los sitios prohibidos en $B+U$.
\paragraph{Cota uniforme.}
Por cálculo espectral, $j_0+\tr T_R^2+\|U\|_{\HS}^2\le C_r q_4$ y $q_4\le C_{M,r}$.
Las entradas de $B,T_s,U$ son $O(k^{-1})$. Por ejemplo
$T_L=B(I-G_R)^{-1}B^\top$ y $U=B(I-G_R)^{-1}B^\top B$;
la cota de cada entrada se obtiene con las normas de las filas y columnas de $B$.
Las diagonales en $j$ son positivas o nulas y
$\|j\|_2^2\le 2j_0\|j\|_\infty=O(q_4/k)$. Para los cubos se multiplica la máxima
entrada absoluta por la suma de cuadrados. El mismo argumento, con $\max|B_{ij}|$,
sirve para la suma de $B_{ij}U_{ij}^2$. Esto prueba~\eqref{eq:q}.
